\documentclass[11pt]{article}
\usepackage[review]{acl}

\usepackage{times}
\usepackage{latexsym}
\usepackage[T1]{fontenc}
\usepackage[utf8]{inputenc}
\usepackage{microtype}
\usepackage{inconsolata}
\usepackage{booktabs}
\usepackage{amsmath}
\usepackage{graphicx}

% Arm names. Defined once so a rename never desynchronises text and tables.
\newcommand{\armB}{\textsc{Base}}
\newcommand{\armA}{\textsc{Ntp}}
\newcommand{\armAA}{\textsc{Aug}}
\newcommand{\armP}{\textsc{Eq}}
\newcommand{\armPF}{\textsc{Code}}
\newcommand{\armC}{\textsc{Ctrl}}
\newcommand{\armMX}{\textsc{Sm-x}}
\newcommand{\armMI}{\textsc{Sm-i}}
\newcommand{\armMP}{\textsc{Pw}}
\newcommand{\armMC}{\textsc{Sm-i+c}}
\newcommand{\armMPC}{\textsc{Pw+c}}

\title{Concept-aware language model training has not been evaluated:\\
a gradient defect, a repair, and a first measurement}

\author{Anonymous submission}

\begin{document}
\maketitle

\section{Methods}

\subsection{Task and objective families}

A causal language model trained with next-token prediction (NTP) maximises the
probability of the single token that was observed at each position. At a
\emph{concept slot}, a position where several words are acceptable, NTP treats
every unobserved acceptable word as an error. Next-concept prediction (NCP)
replaces the target at those positions with a set $C$ of acceptable words and
rewards probability assigned to any member of the set.

We evaluate on SWORDS \citep{lee2021swords}, a lexical substitution benchmark in
which annotators rate candidate substitutes for a target word in context. Ratings
come from between three and ten annotators per candidate. We keep candidates with
acceptability at or above $0.5$ as the acceptable set and candidates at or below
$0.0$ as the rejected set. The labels are human judgements, so no part of the
evaluation depends on model-generated annotation.

Three loss forms appear in this paper. Write $p(c)$ for the model probability of
candidate $c$ at a concept slot, $L_{\text{CLM}}$ for the ordinary causal
language-modelling loss on the same rows, and $C$ for the candidate set.

\paragraph{Per-word.} The mean over candidates of the negative log-probability,
\begin{equation}
L_{\text{pw}} = -\frac{1}{|C|}\sum_{c \in C} \log p(c).
\label{eq:pw}
\end{equation}
Every candidate contributes equally, so the gradient pushes mass onto each member
of the set independently of how the model currently ranks them.

\paragraph{Set marginal.} The negative log of the summed probability,
\begin{equation}
L_{\text{sm}} = -\log \sum_{c \in C} p(c).
\label{eq:sm}
\end{equation}
The gradient of Equation~\ref{eq:sm} is a softmax over the candidates, so it
concentrates on the member the model already prefers. This distinction predicts
that Equation~\ref{eq:pw} spreads probability more evenly than
Equation~\ref{eq:sm}, and we test it directly.

\paragraph{Contrastive.} An InfoNCE term over the human-rejected substitutes,
$\log\sum_{c \in C \cup N} p(c) - \log\sum_{c \in C} p(c)$, where $N$ holds the
$2{,}919$ substitutes SWORDS annotators marked unacceptable. Earlier rounds of
this project mined negatives from WordNet; an audit found that miner produced no
usable wrong-sense negatives, so these human-labelled negatives are the first
valid test of a contrastive term in this setting.

Every trained arm optimises
\begin{equation}
L = w\,L_{\text{CLM}} + \alpha\,L_{\text{concept}} + \beta\,L_{\text{contrast}},
\label{eq:total}
\end{equation}
where $w$ is the retention weight, $\alpha$ scales the concept term and $\beta$
the contrastive term. Reporting $w$ explicitly matters here, because the source
paper's equation and its released code disagree about its value.

\subsection{What the source paper specifies, and what its code does}

\citet{iyer2026beyond} give the NCP objective as
$L(T^*|S) = \frac{1}{|T^*|}\sum_n \log p(t_n \in T^*|S,\Theta)$, maximised. This
is Equation~\ref{eq:pw} at $\alpha = 1$ with no language-modelling term, that is,
$w = 0$. We call this arm \armP{}.

Their released trainer computes
\texttt{loss = base\_loss - (loss\_adjustment / valid\_count)}, where
\texttt{loss\_adjustment / valid\_count} is the mean over slots of the mean over
candidates of $\log p(c)$. That expression is $L_{\text{CLM}} + L_{\text{pw}}$,
which is Equation~\ref{eq:total} at $w = 1$ and $\alpha = 1$. The code therefore
implements a different objective from the printed equation. We call the code's
version \armPF{}.

A third difference separates \armPF{} from what the released code actually
optimises. The candidate log-probabilities are computed inside
\texttt{torch.no\_grad()}, so \texttt{loss\_adjustment} is a constant with
respect to $\Theta$ and $\partial L / \partial \Theta = \partial L_{\text{CLM}} /
\partial \Theta$. The concept term changes the reported loss value and
contributes no gradient. Runs using that trainer were causal language-model
fine-tuning on the concept CSV, and the objective in Equation~\ref{eq:pw} has not
been trained. \armPF{} is its first measurement.

\subsection{Training data views}

Attributing a difference between two arms to their loss requires holding the
training data fixed, so we build four views of the same SWORDS training split and
pad each to the same budget of $1{,}607$ rows.

\textbf{Repeated originals} contains the original sentences, repeated to the
budget, with no candidate sets. \textbf{Augmentation} writes one sentence per
acceptable alternative, which is the data-side procedure the source paper
describes. \textbf{Paper} cycles concept rows to the budget with the observed
word removed from the candidate set. \textbf{Hybrid} mixes concept rows with
repeated originals as replay.

Two design axes cut across these views. A view is \emph{gold-inclusive} when the
observed word stays in the candidate set and \emph{gold-exclusive} when it is
removed; the two give opposite gradient signs on the observed token. A view uses
\emph{replay} when part of its budget holds plain original sentences, which
supplies retention pressure toward the observed continuation. Our proposal
combined gold inclusion with replay; the source paper's view uses neither.

\subsection{Arms}

Table~\ref{tab:arms} lists every arm. \armB{} is the untouched checkpoint and is
included because it is what a practitioner would deploy without any training, and
because a trained arm that loses to it has not earned its compute. \armA{} and
\armAA{} reproduce the source paper's two data-side baselines. \armP{} and
\armPF{} are the printed equation and the repaired code. \armC{} is the
$\alpha = 0$ control that isolates the effect of fine-tuning on this data at all.
\armMX{}, \armMI{}, \armMP{}, \armMC{} and \armMPC{} are ours.

\begin{table}[t]
\centering\footnotesize\setlength{\tabcolsep}{4pt}
\begin{tabular}{lccll}
\toprule
Arm & $w$ & $\alpha$ & Concept loss & Data view \\
\midrule
\armB{}   & n/a & n/a & none          & none \\
\armA{}   & 1   & 0   & none          & repeated \\
\armAA{}  & 1   & 0   & none          & augmentation \\
\armP{}   & 0   & 1   & per-word      & paper, excl. \\
\armPF{}  & 1   & 1   & per-word      & paper, excl. \\
\armC{}   & 1   & 0   & none          & hybrid, incl. \\
\armMX{}  & 1   & 0.5 & set marginal  & hybrid, excl. \\
\armMI{}  & 1   & 0.5 & set marginal  & hybrid, incl. \\
\armMP{}  & 1   & $\alpha$ & per-word & hybrid, incl. \\
\armMC{}  & 1   & 0.5 & set m. $+$ InfoNCE & hybrid, incl. \\
\armMPC{} & 1   & 1   & per-word $+$ InfoNCE & hybrid, incl. \\
\bottomrule
\end{tabular}
\caption{Training arms. $w$ is the weight on the causal language-modelling term
and $\alpha$ scales the concept term (Equation~\ref{eq:total}). \armP{} and
\armPF{} differ only in $w$, so their comparison isolates the term the published
equation omits and the released code includes.}
\label{tab:arms}
\end{table}

\subsection{Metrics}

\textbf{GAP} is the generalised average precision of the model's ranking of
candidates against the annotator ratings, and is the benchmark's official score.
\textbf{p@1} is the fraction of targets whose top-ranked candidate is one
annotators accepted; it is the most directly interpretable number here and moves
in steps of $1/74$ on our evaluation set. \textbf{AUROC} is the separation of
acceptable from rejected substitutes, which is the quantity a contrastive term
claims to improve. \textbf{ALT NLL} is the mean negative log-likelihood of the
acceptable alternatives, so it measures probability mass rather than rank.
\textbf{GOLD NLL} is the negative log-likelihood of the word that actually
appeared, and it is the cost side of the trade: an objective can improve every
ranking metric by moving mass off the observed word. \textbf{NTP PPL} is ordinary
perplexity on held-out text and reports whether the model remains a language
model. AUROC and ALT NLL come apart in a way that matters, since a model can rank
perfectly while holding most of its mass on rejected candidates.

\subsection{Protocol}

We train Llama-3.2-1B for two epochs at learning rate $1\times10^{-5}$ with a
cosine schedule and 10\% warmup, on a single T4, with three seeds
($42$, $7$, $123$). Every arm sees the same $1{,}607$ rows, the same schedule and
the same seeds, so any difference between two arms comes from the fields
Table~\ref{tab:arms} varies. We evaluate on $74$ held-out SWORDS targets that
never appear in training.

Comparisons use a paired bootstrap over targets with $5{,}000$ resamples,
computed per seed. We fixed the claim rule before reading any result: a claim is
permitted only when all three seeds agree in sign \emph{and} every seed's
confidence interval excludes zero. Tables report the mean over seeds with the
standard deviation across seeds. The gate arms were fixed in advance so that arms
added later could not move the pass criterion.

\section{Results}

\subsection{Repairing the gradient produces the strongest arm}

Table~\ref{tab:main} gives the main result. \armP{}, the published equation
implemented literally, ranks substitutes well and damages the model: GOLD NLL
rises to $6.159$ and perplexity to $9.807$, both worse than the untouched
checkpoint at $4.446$ and $7.846$. Adding the language-modelling term that the
released code contains, and repairing the gradient, gives \armPF{}. It holds
\armP{}'s ranking quality within noise (GAP $0.4553$ against $0.4559$, AUROC
$0.7131$ against $0.7123$, p@1 $0.4369$ against $0.4324$) and recovers most of
the loss: GOLD NLL falls to $4.966$ and perplexity to $8.739$.

The collapse under \armP{} follows from the missing retention term. Any
description of the published equation as damaging to the language model describes
the equation as printed and not the method the authors ran.

\begin{table*}[t]
\centering\small\setlength{\tabcolsep}{4pt}
\begin{tabular}{lcccccc}
\toprule
Arm & GAP $\uparrow$ & p@1 $\uparrow$ & AUROC $\uparrow$ & ALT NLL $\downarrow$ & GOLD NLL $\downarrow$ & NTP PPL $\downarrow$ \\
\midrule
\armB{} untouched base & 0.4009 & 0.2432 & 0.6711 & 8.049 & 4.446 & 7.846 \\
\armA{} repeated originals & 0.3961\,{\tiny$\pm$0.0008} & 0.2973\,{\tiny$\pm$0.0000} & 0.6604\,{\tiny$\pm$0.0004} & 9.520\,{\tiny$\pm$0.011} & 4.527\,{\tiny$\pm$0.005} & 9.528\,{\tiny$\pm$0.019} \\
\armAA{} NCP augmentation & 0.4111\,{\tiny$\pm$0.0011} & 0.3063\,{\tiny$\pm$0.0078} & 0.6792\,{\tiny$\pm$0.0005} & 7.996\,{\tiny$\pm$0.029} & 4.516\,{\tiny$\pm$0.005} & 9.040\,{\tiny$\pm$0.009} \\
\armP{} published equation & 0.4559\,{\tiny$\pm$0.0008} & 0.4324\,{\tiny$\pm$0.0135} & 0.7123\,{\tiny$\pm$0.0005} & 5.997\,{\tiny$\pm$0.007} & 6.159\,{\tiny$\pm$0.016} & 9.807\,{\tiny$\pm$0.108} \\
\armPF{} released code, repaired & 0.4553\,{\tiny$\pm$0.0004} & 0.4369\,{\tiny$\pm$0.0078} & 0.7131\,{\tiny$\pm$0.0011} & 6.366\,{\tiny$\pm$0.003} & 4.966\,{\tiny$\pm$0.016} & 8.739\,{\tiny$\pm$0.051} \\
\armMI{} set marginal & 0.4016\,{\tiny$\pm$0.0006} & 0.3018\,{\tiny$\pm$0.0078} & 0.6712\,{\tiny$\pm$0.0005} & 8.575\,{\tiny$\pm$0.023} & 4.167\,{\tiny$\pm$0.015} & 8.372\,{\tiny$\pm$0.023} \\
\armMP{} per-word, $\alpha{=}1$ & 0.4121\,{\tiny$\pm$0.0027} & 0.2838\,{\tiny$\pm$0.0135} & 0.6876\,{\tiny$\pm$0.0024} & 7.440\,{\tiny$\pm$0.055} & 4.206\,{\tiny$\pm$0.005} & 8.201\,{\tiny$\pm$0.024} \\
\armMP{} per-word, $\alpha{=}8$ & 0.4319\,{\tiny$\pm$0.0049} & 0.3423\,{\tiny$\pm$0.0207} & 0.7024\,{\tiny$\pm$0.0026} & 6.568\,{\tiny$\pm$0.084} & 4.530\,{\tiny$\pm$0.040} & 8.283\,{\tiny$\pm$0.033} \\
\bottomrule
\end{tabular}
\caption{SWORDS results on 74 held-out targets, mean over three seeds with
standard deviation across seeds. \armB{} is deterministic and has one value.
\armPF{} leads on every ranking metric. Our arms lead on GOLD NLL and perplexity,
which is the retention side of the trade rather than the benchmark's target.}
\label{tab:main}
\end{table*}

\subsection{Every arm of ours loses to the repaired objective}

Table~\ref{tab:vsp1f} compares each of our arms against \armPF{} under the
pre-registered rule. All six lose on p@1 and on AUROC, with all seeds agreeing
and confidence intervals excluding zero in eleven of the twelve cells. All six
win on GOLD NLL by between $0.76$ and $0.80$.

The direction is consistent across our whole design. Our arms preserve the
observed word better and rank alternatives worse. Both effects follow from the
same two choices: gold inclusion keeps the observed word in the candidate set,
and replay spends part of the budget on plain original sentences. Both supply
retention pressure, so they buy GOLD NLL and cost ranking. On a benchmark scored
by ranking, that places our design on the wrong side of the trade.

\begin{table}[t]
\centering\footnotesize\setlength{\tabcolsep}{4pt}
\begin{tabular}{lccc}
\toprule
Arm vs \armPF{} & p@1 & AUROC & GOLD NLL \\
\midrule
\armC{} & $-0.1261$ & $-0.0429$ & $-0.7873$ \\
\armMX{} & $-0.1216^{\dagger}$ & $-0.0344$ & $-0.7965$ \\
\armMI{} & $-0.1351$ & $-0.0418$ & $-0.7994$ \\
\armMP{} & $-0.1532$ & $-0.0255$ & $-0.7604$ \\
\armMC{} & $-0.1306$ & $-0.0414$ & $-0.7993$ \\
\armMPC{} & $-0.1532$ & $-0.0266$ & $-0.7584$ \\
\bottomrule
\end{tabular}
\caption{Paired bootstrap deltas against \armPF{}, mean over three seeds.
Negative is worse for p@1 and AUROC and better for GOLD NLL. All entries have all
three seeds agreeing in sign; $\dagger$ marks the one cell where a confidence
interval includes zero. Every other entry satisfies the pre-registered claim
rule. Arm definitions are in Table~\ref{tab:arms}.}
\label{tab:vsp1f}
\end{table}

\subsection{Raising the concept weight does not close the gap}

The concept weight $\alpha$ was inherited from the source paper's equation and
had never been tuned, so a reasonable objection to Table~\ref{tab:vsp1f} is that
our arms are simply under-weighted. Table~\ref{tab:alpha} sweeps it.

GAP rises with $\alpha$ and saturates. The gain per doubling falls from $+0.0070$
to $+0.0013$, leaving a ceiling near $0.432$ against \armPF{}'s $0.4553$. GOLD
NLL rises throughout and passes the untouched checkpoint's $4.446$ at
$\alpha = 6$, so beyond that point the arm pays more than an untrained model on
the observed word while still ranking below \armPF{}. Reaching \armPF{}'s GOLD
NLL of $4.966$ by extrapolation requires roughly $\alpha = 16$, at a projected
GAP near $0.433$. Our arms do not appear to dominate even at matched retention.

\begin{table}[t]
\centering\footnotesize\setlength{\tabcolsep}{4pt}
\begin{tabular}{lccccc}
\toprule
$\alpha$ & GAP & $\Delta$GAP & AUROC & ALT NLL & GOLD NLL \\
\midrule
1 & 0.4121 &  & 0.6876 & 7.440 & 4.206 \\
2 & 0.4191 & $+0.0070$ & 0.6932 & 7.048 & 4.279 \\
4 & 0.4270 & $+0.0079$ & 0.6988 & 6.744 & 4.405 \\
6 & 0.4306 & $+0.0036$ & 0.7011 & 6.628 & 4.480 \\
8 & 0.4319 & $+0.0013$ & 0.7024 & 6.568 & 4.530 \\
\bottomrule
\end{tabular}
\caption{Concept weight sweep on \armMP{}. $\Delta$GAP is the change from the
previous row. Gains saturate while GOLD NLL keeps rising, crossing the untouched
checkpoint ($4.446$) at $\alpha = 6$.}
\label{tab:alpha}
\end{table}

\subsection{Two design choices do not survive}

Table~\ref{tab:ablation} reports the two ablations that carried predictions.

We registered the hypothesis that gold-\emph{inclusive} candidate sets would beat
gold-exclusive ones, on the reasoning that keeping the observed word prevents the
model from unlearning it. The opposite holds in all three seeds: \armMX{} beats
\armMI{} on GAP ($0.4065$ against $0.4016$), AUROC ($0.6787$ against $0.6712$)
and ALT NLL ($8.015$ against $8.575$), with non-overlapping seed ranges. \armPF{}
is also gold-exclusive. We report this as a descriptive result, since the paired
bootstrap for this pair was not computed in the present run.

The contrastive term is null on both base objectives. Adding InfoNCE over
human-labelled rejected substitutes moves AUROC by $+0.0004$ on the set marginal,
with the sign disagreeing across seeds, and leaves the per-word arm unchanged
within noise. These are the hardest negatives available in this setting, so the
result is not explained by negative quality. Pretrained models already separate
acceptable from rejected substitutes at AUROC $0.671$ before any training, which
suggests little remains for a contrastive term to learn at this data scale.

\begin{table}[t]
\centering\footnotesize\setlength{\tabcolsep}{4pt}
\begin{tabular}{lccc}
\toprule
Arm & GAP $\uparrow$ & AUROC $\uparrow$ & ALT NLL $\downarrow$ \\
\midrule
gold-exclusive \armMX{} & 0.4065 & 0.6787 & 8.015 \\
gold-inclusive \armMI{} & 0.4016 & 0.6712 & 8.575 \\
\armMI{} $+$ InfoNCE & 0.4020 & 0.6716 & 8.568 \\
\armMP{} $+$ InfoNCE & 0.4112 & 0.6865 & 7.452 \\
\armMP{} alone & 0.4121 & 0.6876 & 7.440 \\
\bottomrule
\end{tabular}
\caption{Design ablations, mean over three seeds. Top: gold exclusion beats gold
inclusion at matched objective and data. Bottom: adding InfoNCE over
human-rejected substitutes changes neither base objective beyond seed noise.}
\label{tab:ablation}
\end{table}

\subsection{What the data view contributes}

\armMP{} beats \armMI{} on GAP ($+0.0105$), AUROC ($+0.0164$) and ALT NLL
($-1.1344$), all under the claim rule. The per-word form spreads probability
across the candidate set where the set marginal concentrates it on the member the
model already prefers, which is what Equations~\ref{eq:pw} and \ref{eq:sm}
predict. This is the one component of our design that behaves as intended.

The remaining gap between \armMP{} and \armPF{} spans two fields at once, the
candidate set and the data view, so the present run cannot divide it. We report
that as an open attribution rather than assigning it to either.

\section*{Limitations}

We train a single 1B model. The source paper used models in the 7B to 8B range,
and the effect of a concept objective may depend on scale, so our comparison is
internally matched but externally narrow. The training set holds $1{,}607$ rows
and the evaluation $74$ targets, which quantises p@1 at $1/74$ and limits the
resolution of every ranking metric. Selection of $\alpha$ used the development
set that Tables~\ref{tab:main} and \ref{tab:alpha} report, so the tuned rows are
optimistic and a held-out test measurement is still owed. The gold-exclusive
result in Table~\ref{tab:ablation} rests on seed ranges rather than a paired
bootstrap. Finally, our conclusion about the released implementation applies to
the code in that repository; we cannot rule out that the published numbers came
from a version that was not released.

\bibliography{custom}

\end{document}
